% DRAFT, rewrite
\section{Introduction}

\subsection{Purpose}
PoliVerse is an unofficial app for students of the Politecnico di Milano. It shows
courses, materials, the timetable, the career and the campus in one place. This
document specifies one update to it: how the app loads data and how it shows data
that is still on its way.

Today several screens wait for every request to answer before showing anything, the
app cannot tell ``empty'' from ``not here yet'', and nothing is prepared before the
student asks for it. The update serves the following goals.

\goal{G1}{A student who opens a screen sees its first useful content without waiting for everything on that screen.}
\goal{G2}{A student is never shown a statement or a number that is false only because data is still on its way.}
\goal{G3}{A student who is about to attend a lecture finds what they need for it already on the phone.}
\goal{G4}{A student can consult the timetable, the campus, the rooms and their courses with no network connection.}
\goal{G5}{A student notices what changed since they last looked without reading everything again.}
\goal{G6}{Work the student did not ask for never spends their mobile data or battery beyond what they allowed, and stays within a fixed number of requests to the Politecnico.}
\goal{G7}{The developer can tell from measurements taken on a phone whether a change to loading helped.}

\subsection{Scope}
The machine is the PoliVerse app on iPhone, with its widgets and background tasks.
The world is the student, the phone and its network, and the Politecnico's servers.
The project runs no server of its own.

\begin{center}\small
\begin{tabularx}{\linewidth}{@{}X l@{}}
  \toprule
  Phenomenon & Kind \\
  \midrule
  A lecture takes place in a room at a given time & World \\
  A room is booked or a booking is withdrawn & World \\
  A grade is published, an exam sitting opens, a lecture is moved & World \\
  The student travels through a place with no signal & World \\
  The student opens a screen, scrolls, taps a retry control & Shared, world controlled \\
  A server answers a request, fails, or does not answer & Shared, world controlled \\
  The network becomes available, metered or unavailable & Shared, world controlled \\
  iOS grants background time; reports Low Data Mode, Low Power Mode and heat & Shared, world controlled \\
  The app sends a request to a server & Shared, machine controlled \\
  The app shows, updates or animates a section & Shared, machine controlled \\
  The app delivers a local notification or a haptic & Shared, machine controlled \\
  The app keeps a stored copy, orders requests, records the ledger & Machine \\
  \bottomrule
\end{tabularx}
\end{center}

In scope: showing sections as their parts arrive, the order of requests, storing and
ageing parts one by one, deadlines for slow parts, prefetching from the timetable, an
offline pack, refreshing at booking boundaries, noticing changes, the limits on
speculative requests, and the loading ledger.

Out of scope: any server operated by the project and therefore real push; predicting
the next screen from the student's habits; changes to what the screens contain;
loading of recordings' video, which has its own download manager.

\subsection{Definitions, acronyms, abbreviations}
\printterms

\noindent Acronyms: RASD, Requirements Analysis and Specification Document; DD, Design
Document; UC, use case; HIG, Human Interface Guidelines; TTL, time to live.

\noindent A number followed by \initial{} is an initial value. It is a requirement as
written, and it is confirmed or changed once the loading ledger (\idref{R34}) has
produced a baseline.

\subsection{Revision history}
See the table under the title.

\subsection{Reference documents}
Listed in section 6.

\subsection{Document structure}
Section 2 describes the update from the outside: scenarios, the domain, the
functions, the users and the assumptions about the world. Section 3 lists the
requirements: interfaces, functional requirements with their use cases, performance
requirements, constraints and system attributes. Section 4 states what is checked
formally. Sections 5 and 6 hold the effort spent and the references.
